\BourbakiExercisesSection

\begin{enumerate}
\def\labelenumi{\arabic{enumi}.}
\item
  证明：若 E 是一个含无穷多个元素的域 K 上的向量空间 \(\neq\{0\}\) ，则 E 的生成系所成的集对于序关系 \(\supset\) 不是归纳的（作出 E 的生成系的一个递减序列 \((S_{n})\) ，使得诸 \(S_{n}\) 的交是空的）。
\item
  设 K 为一个域，L 为 K 的子域，I 为一个集合，且 \((x_{i})_{1\leqslant i\leqslant m}\) 为 \(K_{s}^{I}\) 中的一个自由向量族，使得每个 \(x_{i}\) 的坐标都属于 L。设 V 为 \(K_{s}^{I}\) 中由这些 \(x_{i}\) 生成的向量子空间；证明存在 I 的一个含 m 个元素的有限子集 J，使得对 V 的每个向量 \(z=(\zeta_{\alpha})_{\alpha\in\mathbb{I}}\) ，所有 \(\zeta_{\alpha}\) 都属于由 L 与 m 个元素 \(\zeta_{\beta}\) （其中 \(\beta\in J\) ）生成的 K 的子域。（通过将第 5 节的推论 3 应用于空间 \(L_{s}^{I}\) ，证明可化归于存在 m 个指标 \(\beta_{i}\in I\) （ \(1\leqslant i\leqslant m\) ）使得 \(\Pr_{\beta_{i}}(x_{j})=\delta_{ij}\) 的情形。）
\item
  \begin{enumerate}
  \def\labelenumii{(\alph{enumii})}
  \tightlist
  \item
    设 G 是一个群，M 是 G 的一个无限子集。证明：若 \(G'\) 是由 M 生成的 G 的子群，则 \(\text{Card}(G') = \text{Card}(M)\) 。
  \end{enumerate}
\end{enumerate}

\begin{enumerate}
\def\labelenumi{(\alph{enumi})}
\setcounter{enumi}{1}
\item
  设 \(\mathbf{K}\) 为一个域， \(\mathbf{M}\) 为 \(\mathbf{K}\) 的一个无限子集。证明：若 \(\mathbf{K}'\) 是 \(\mathbf{K}\) 的由 \(\mathbf{M}\) 生成的子域，则 \(\operatorname{Card}(\mathbf{K}') = \operatorname{Card}(\mathbf{M})\) 。（考虑子集的两个序列 \((\mathbf{A}_n)_{n \geqslant 0}, (\mathbf{P}_n)_{n \geqslant 0}\) （这些子集属于 \(\mathbf{K}\) ），使得 \(\mathbf{A}_0 = \mathbf{M}\) ， \(\mathbf{P}_n\) 是 \(\mathbf{K}^*\) 的由 \(\mathbf{A}_n \cap \mathbf{K}^*\) 生成的乘法子群，而 \(\mathbf{A}_{n+1}\) 是 \(\mathbf{K}\) 的由 \(\mathbf{P}_n\) 生成的加法子群；然后应用 (a)。）
\item
  设 K 为一个域，I 为一个无限集；证明 \(\operatorname{Card}(\mathbf{K}) \leqslant \dim(\mathbf{K}_{s}^{\mathbf{I}})\) 。（将其归结为 I = N 的情形，并用反证法论证。设 B 为 \(K_{s}^{N}\) 的满足 \(\operatorname{Card}(\mathbf{B}) < \operatorname{Card}(\mathbf{K})\) 的一组基，并设 L 为 K 中由 B 的所有元素的坐标生成的子域；则 \(\operatorname{Card}(\mathbf{L}) < \operatorname{Card}(\mathbf{K})\) 。构造 K 中元素的一个序列 \((\xi_{n})_{n \geqslant 1}\) ，使得对所有 n， \(\xi_{n}\) 不属于 K 中由 L 与诸 \(\xi_{i}\) （其指标 i \textless{} n ）生成的子域；然后应用习题 2，通过考虑点 \(x = (\xi_{n})\) （属于 \(K_{s}^{N}\) ）得出矛盾。）
\item
  由(c)可推出，对任意域K及任意无限集I，
\end{enumerate}

\[
\dim (\mathrm{K} _ {s} ^ {\mathrm{I}}) = (\operatorname{Card} (\mathrm{K})) ^ {\operatorname{Card} (\mathrm{I})}
\]

(``Erdös-Kaplansky 定理'').

4 给出一个无限序列 \((\mathbf{F}_n)\) 的例子，该序列由某个向量空间 \(\mathbf{E}\) 的向量子空间组成，使得 \(\operatorname{codim}\left(\bigcap_{n} \mathbf{F}_n\right) > \sum_{n} \operatorname{codim}(\mathbf{F}_n)\) 。（取 \(\operatorname{codim}(\mathbf{F}_n) = 1\) 对所有 \(n\) 成立，并利用习题 3 \((d)\) 。）

\begin{enumerate}
\def\labelenumi{\arabic{enumi}.}
\setcounter{enumi}{4}
\tightlist
\item
  设 \((\mathbf{H}_{\lambda})_{\lambda \in L}\) 是向量空间 \(\mathbf{E}\) （在域 \(\mathbf{K}\) 上）的一族超平面（经过 0），它构成 \(\mathbf{E}\) 的一个覆盖。
\end{enumerate}

\begin{enumerate}
\def\labelenumi{(\alph{enumi})}
\item
  证明：若 K 有限，则 \(\operatorname{Card}(\mathbf{L}) \geqslant 1 + \operatorname{Card}(\mathbf{K})\) ；若 K 无限，则 \(\operatorname{Card}(\mathbf{L}) \geqslant \aleph_{0} = \operatorname{Card}(\mathbf{N})\) 。（对 r 作归纳证明：若 K 至少有 r 个元素，则 E 不能是 r 个超平面的并。）用例子说明，在没有附加假设的情况下，这些不等式不能改进（当 K 无限时，考虑空间 \(\mathbf{E} = \mathbf{K}_{s}^{(\mathbf{N})}\) ）。
\item
  若 E 是有限维的，证明 \(\operatorname{Card}(\mathbf{L}) \geqslant 1 + \operatorname{Card}(\mathbf{K})\) 。（对 \(\dim(\mathbf{E})\) 作归纳论证。）这一结论在无限维时不再成立；考虑 \(K^{(N)}\)
\end{enumerate}

¶ 6. 设 S 是向量 K-空间 E 的一个非空有限子集，且不包含 0。假设存在一个整数 \(k \geqslant 1\) 具有如下性质：(*) 对所有 \(x \in S\) ，存在 \(S = \{x\}\) 到 \(h\) 个自由子集的一个划分，其中 \(h \leqslant k\) 。

对于这样的划分 \((\mathbf{N}_i)_{1\leqslant i\leqslant h}\) 与每个整数序列 \((r_j)_{1\leqslant j\leqslant m}\) （整数取自 \(\{1,h\}\) ），子空间 \(\mathrm{F}_{r_1r_2\dots r_m}\) （属于 \(\mathbf{E}\) ）按对 \(m\) 的归纳定义如下： \(\mathrm{F}_{r_1}\) 是由 \(\mathrm{N}_{r_1},\mathrm{F}_{r_1\dots r_p}\) 生成的子空间，后者为由 \(\mathrm{F}_{r_1\dots r_{p-1}}\) 与 \(\mathrm{N}_{r_p}\) 的交所生成的子空间。假设：

(**) 对所有 \(x \in S\) 以及每个划分 \((N_i)_{1 \leqslant i \leqslant h}\) （将 \(S = \{x\}\) 分为至多 \(k\) 个自由子集），至少存在一个序列 \((r_j)_{1 \leqslant j \leqslant m}\) ，使得

\[
x \notin \mathrm{F} _ {r _ {1} r _ {2} \dots r _ {m}}.
\]

\begin{enumerate}
\def\labelenumi{(\alph{enumi})}
\item
  设 \(n\) 是整数 \(m \geqslant 1\) （对 \(x \in S\) 的所有可能选取）中的最小者，并设 \((\mathbf{N}_i)_{1 \leqslant i \leqslant h}\) 和 \((r_j)_{1 \leqslant j \leqslant m}\) 满足条件 (*) 和 (**)。证明必然有 \(n = 1\) 。（用反证法论证，假设 \(x \notin F_{r_1 r_2 \ldots r_n}\) 且 \(n > 1\) 。则必有 \(x \in F_{r_n}\) 。为简单起见，记 \(E_p = F_{r_1 r_2 \ldots r_p}\) （ \(1 \leqslant p \leqslant n\) ），并写 \(x = \sum_i \alpha_i y_i\) ，其中 \(\alpha_i \in K\) ， \(y_i \in N_{r_n}\) ；至少有一个 \(y_i\) 不属于 \(E_{n-1}\) ，把它记作 \(y\) ；然后令 \(N = N_{r_n} - \{y\}\) ，令 \(N_i' = N_i\) （ \(1 \leqslant i \leqslant h\) 且 \(i \neq r_n\) ），令 \(N_{r_n}' = N \cup \{x\}\) ，使得 \((\mathbf{N}_i')_{1 \leqslant i \leqslant h}\) 是由自由子集组成的 \(S - \{y\}\) 的一个划分。定义 \(E_0' = E\) ，并对 \(1 \leqslant p \leqslant n\) 把 \(E_p'\) 定义为由 \(E_{p-1}'\) 与 \(N_{r_p}'\) 的交生成的子空间。现在 \(y \notin E_{n-1}\) ；令 \(q\) 为使 \(y \notin E_q\) 的最小整数；证明 \(E_q' \not\subset E_q\) 。令 \(s\) 为使 \(E_s' \not\subset E_s\) 的最小整数；证明 \(r_s = r_n\) ；最后由关系 \(y \in E_s, E_{s-1}' \subset E_{s-1}\) 与 \(E_s' \not\subset E_s\) 推出矛盾。
\item
  由 (a) 推出，在给定假设下，存在 S 到 h 个自由子集的一个划分，其中整数 \(h \leqslant k\) 。
\end{enumerate}

\begin{enumerate}
\def\labelenumi{\arabic{enumi}.}
\setcounter{enumi}{6}
\tightlist
\item
  设 S 是向量 K-空间 E 的一个非空有限子集，且不包含 0。假设存在一个整数 \(k \geqslant 1\) 使得对于每个子集
\end{enumerate}

S 的每个子集 T，有 \(\text{Card}(T) \leqslant k.\text{rg}(T)\) 。证明存在将 S 分为 h 个自由子集的一个划分，其中某个整数 \(h \leqslant k\) 。（证明 S 满足习题 6 的条件（ \(^{**}\) ），对 \(\text{Card}(S)\) 作归纳论证，并在所有子空间 \(F_{r_{1}\ldots r_{m}}\) 中取一个维数尽可能小的子空间；若 G 是这样的子空间，且 j 是使 \(\text{Card}(G \cap N_{j})\) 为最小的指标，则证明

\[
\begin{array}{r l} k. \operatorname{Card} (\mathrm{G} \cap \mathrm{N} _ {j}) & \leqslant \operatorname{Card} (\mathrm{G} \cap (\mathrm{S} - \{x \})) \\ & \leqslant \operatorname{Card} (\mathrm{G} \cap \mathrm{S}) \leqslant k. \operatorname{Card} (\mathrm{G} \cap \mathrm{N} _ {j}).) \end{array}
\]

\begin{enumerate}
\def\labelenumi{\arabic{enumi}.}
\setcounter{enumi}{7}
\item
  设 E 是一个向量空间。证明 E 的每个在 End(E) 中不是 0 的右因子的自同态都是满射的（§ 2，习题 7）。
\item
  设K是一个域。
\end{enumerate}

\begin{enumerate}
\def\labelenumi{(\alph{enumi})}
\item
  在向量空间 \(\mathbf{E} = \mathrm{K}_s^{(\mathbf{Z})}\) 中，设 \((e_n)_{n\in \mathbf{Z}}\) 为典范基；设 \(v\) 为 \(\mathbf{E}\) 的由关系 \(v(e_n) = e_{n + 1}\) （对所有 \(n\in \mathbf{Z}\) ）定义的自同构。若 \(u = 1_{\mathbb{E}} - v\) ，证明 \(u\) 是 \(\mathbf{E}\) 的一个单射自同态，但 \(\dim (\operatorname {Coker}(u)) = 1\) 。
\item
  在向量空间 \(\mathbf{E} = \mathrm{K}_s^{(\mathbb{N})}\) 中，设 \((e_n)_{n\in \mathbb{N}}\) 为典范基；设 \(v\) 为 \(\mathbf{E}\) 的由 \(v(e_0) = e_0\) ， \(v(e_n) = e_{n - 1} + e_n\) （对 \(n\geqslant 1\) ）定义的自同态。证明 \(v\) 是 \(\mathbf{E}\) 的一个自同构，且若 \(u = 1_{\mathbb{E}} - v\) ，则 \(u\) 是 \(\mathbf{E}\) 的满足 \(\dim (\mathrm{Ker}(u)) = 1\) 的一个满射自同态。
\item
  设 I 为一个非空集合。证明向量空间 \(E = K_{s}^{(I)}\) 的每个自同构 u 都可以写成两个自同构之差 u = v - w，除非 K 只有两个元素且 I 仅由一个元素组成。（注意只需证明某个特定的自同构 u 具有所述形式即可；当 K 是两个元素的域时，分别考虑 \(\text{Card}(I) = 2\) 与 \(\text{Card}(I) = 3\) 的情形。）
\end{enumerate}

设 E 是域 K 上的向量空间。证明 E 的每个自同态 u 要么是自同构，要么可以写成 u = v - w，其中 v 和 w 是 E 的自同构。（注意，u 总可以换成 \(s_{1} \circ u \circ s_{2}\) ，其中 \(s_{1}\) 和 \(s_{2}\) 是 E 的两个自同构。区分两种情形，视 \(\dim(\text{Ker}(u)) \leqslant \dim(\text{Coker}(u))\) 还是 \(\dim(\text{Ker}(u)) > \dim(\text{Coker}(u))\) 而定；当 \(\text{rg}(u)\) 有限时，第一种情形总成立。当

\[
\dim (\operatorname{Ker} (u)) \leqslant \dim (\operatorname{Coker} (u)),
\]

它可以简化为以下情形： \(\operatorname{Ker}(u) \cap \operatorname{Im}(u) = \{0\}\) ；若

\[
\dim (\operatorname{Ker} (u)) = \dim (\operatorname{Coker} (u)),
\]

还可以假设 \(\operatorname{Im}(u) + \operatorname{Ker}(u) = \mathbf{E}\) ，然后应用习题 9 (c)。若 \(\dim (\operatorname {Ker}(u)) < \dim (\operatorname {Coker}(u))\) ，则在 \(\mathbf{E}\) 中存在 \(\operatorname {Ker}(u)\) 的一个形如 \(\mathbf{W} \oplus \operatorname{Im}(u)\) 的补子空间，满足 \(\dim (\mathbf{W}) = \dim (\operatorname {Im}(u)) = \dim (\mathbf{E})\) ；应用开头的注记并取 \(s_1(u(\operatorname {Im}(u))) = \operatorname{Im}(u)\) ，并利用习题 9 (a) 和 (c)。当 \(\dim (\operatorname {Ker}(u)) > \dim (\operatorname {Coker}(u))\) 时，可化归于 \(\operatorname{Ker}(u) \subset \operatorname{Im}(u)\) 的情形；这次取 \(s_1(\operatorname{Ker}(u)) = u(\operatorname{Ker}(u))\) ，并利用前面诸情形及习题 9 的结果。）

\begin{enumerate}
\def\labelenumi{\arabic{enumi}.}
\setcounter{enumi}{10}
\tightlist
\item
  设 E、F 是域 K 上的两个向量空间， \(u: E \rightarrow F\) 是一个线性映射。证明对于 F 的每个向量子空间 V，
\end{enumerate}

\[
\dim (u ^ {- 1} (\mathbf {V})) = \dim (\mathbf {V} \cap \operatorname{Im} (u)) + \dim (\operatorname{Ker} (u)).
\]

\begin{enumerate}
\def\labelenumi{\arabic{enumi}.}
\setcounter{enumi}{11}
\tightlist
\item
  设 E、F 为两个向量空间，u、v 为从 E 到 F 的两个线性映射。
\end{enumerate}

\begin{enumerate}
\def\labelenumi{(\alph{enumi})}
\tightlist
\item
  证明以下不等式成立：
\end{enumerate}

\[
\begin{array}{c} \operatorname{rg} (v) \leqslant \operatorname{rg} (u + v) + \operatorname{rg} (u) \\ \operatorname{rg} (u + v) \leqslant \inf (\dim (\mathrm{E}), \dim (\mathrm{F}), \operatorname{rg} (u) + \operatorname{rg} (v)). \end{array}
\]

若E和F是有限维的，证明 \(\operatorname{rg}(u + v)\) 可以取满足不等式的任意整数值

\[
\left| \operatorname{rg} (u) - \operatorname{rg} (v) \right| \leqslant \operatorname{rg} (u + v) \leqslant \inf (\dim (\mathrm{E}), \dim (\mathrm{F}), \operatorname{rg} (u) + \operatorname{rg} (v)).
\]

\begin{enumerate}
\def\labelenumi{(\alph{enumi})}
\setcounter{enumi}{1}
\tightlist
\item
  证明
\end{enumerate}

\[
\dim (\operatorname{Ker} (u + v)) \leqslant \dim (\operatorname{Ker} (u) \cap \operatorname{Ker} (v)) + \dim (\operatorname{Im} (u) \cap \operatorname{Im} (v)).
\]

\begin{enumerate}
\def\labelenumi{(\alph{enumi})}
\setcounter{enumi}{2}
\tightlist
\item
  证明
\end{enumerate}

\[
\dim (\operatorname{Coker} (u + v)) \leqslant \dim (\operatorname{Coker} (u)) + \operatorname{rg} (v).
\]

（若 \(\mathbf{V}\) 是 \(\operatorname{Ker}(v)\) 在 \(\mathbf{E}\) 中的一个补子空间，注意

\[
u (\operatorname{Ker} (v)) \subset \operatorname{Im} (u + v) \quad \text { and } \quad \dim (u (\mathbf {V})) \leqslant \operatorname{rg} (v).)
\]

\begin{enumerate}
\def\labelenumi{\arabic{enumi}.}
\setcounter{enumi}{12}
\tightlist
\item
  设 E、F、G 为三个向量空间，且 \(u: \mathbf{E} \to \mathbf{F}\) , \(v: \mathbf{F} \to \mathbf{G}\) 为两个线性映射。
\end{enumerate}

\begin{enumerate}
\def\labelenumi{(\alph{enumi})}
\tightlist
\item
  证明存在将 E 分解为 \(\operatorname{Ker}(u)\) 与两个子空间 M, N 的直和，使得 \(\operatorname{Ker}(v \circ u) = \mathbf{M} \oplus \operatorname{Ker}(u)\) 和
\end{enumerate}

\[
\operatorname{Im} (v \circ u) = v (u (\mathbf {N})).
\]

\begin{enumerate}
\def\labelenumi{(\alph{enumi})}
\setcounter{enumi}{1}
\item
  \(\dim(\operatorname{Ker}(u)) \leqslant \dim(\operatorname{Ker}(v \circ u)) \leqslant \dim(\operatorname{Ker}(u)) + \dim(\operatorname{Ker}(v))\) ；若 \(\operatorname{Ker}(u)\) 和 \(\operatorname{Ker}(v)\) 是有限维的，证明 \(\dim(\operatorname{Ker}(v \circ u))\) 可以取满足上述不等式的每一个整数值。
\item
  以下等式成立：
\end{enumerate}

\[
\begin{array}{l} \operatorname{rg} (u) = \operatorname{rg} (v \circ u) + \dim (\operatorname{Im} (u) \cap \operatorname{Ker} (v)) \\ \operatorname{rg} (v) = \operatorname{rg} (v \circ u) + \operatorname{codim} _ {\mathbb {F}} (\operatorname{Im} (u) + \operatorname{Ker} (v)). \end{array}
\]

若 F 是有限维的，证明 \(\operatorname{rg}(v \circ u)\) 可以取满足该不等式的任意整数值

\[
\sup (0, \operatorname{rg} (u) + \operatorname{rg} (v) - n) \leqslant \operatorname{rg} (v \circ u) \leqslant \inf (\operatorname{rg} (u), \operatorname{rg} (v)).
\]

\begin{enumerate}
\def\labelenumi{\arabic{enumi}.}
\setcounter{enumi}{13}
\item
  设 E、F、G 为三个向量空间，且 \(u: \mathbf{E} \to \mathbf{F}\) , \(w: \mathbf{E} \to \mathbf{G}\) 两个线性映射。证明若 \(u(0) \subset w(0)\) ，存在一个线性映射 \(v\) ，使得 \(w = v \circ u\) （参见第2节，习题8）。
\item
  设E，F是域K上的两个向量空间， \(u: E \rightarrow F\) 是一个线性映射。
\end{enumerate}

\begin{enumerate}
\def\labelenumi{(\alph{enumi})}
\item
  下列条件等价：（1） \(\operatorname{Ker}(u)\) 是有限维的；（2）存在一个线性映射 \(v\) （由 \(\mathbf{F}\) 到 \(\mathbf{E}\) ），使得 \(v \circ u = 1_{\mathbf{E}} + w\) ，其中 \(w\) 是 \(\mathbf{E}\) 的一个有限秩自同态；（3）对每个向量空间 \(\mathbf{G}\) （在 \(\mathbf{K}\) 上）与每个线性映射 \(f: \mathbf{G} \to \mathbf{E}\) ，关系 \(\operatorname{rg}(u \circ f) < +\infty\) 蕴含 \(\operatorname{rg}(f) < +\infty\) 。
\item
  以下条件等价：（1）Coker(u)是有限维的；（2）存在一个线性映射v：F→E，使得 \(u \circ v = 1_{F} + w\) ，其中w是F的一个有限秩自同态；（3）对于K上的每个向量空间G和每个线性映射 \(g: F \to G\) ，关系 \(\operatorname{rg}(g \circ u) < +\infty\) 蕴含 \(\operatorname{rg}(g) < +\infty\) .
\end{enumerate}

\begin{enumerate}
\def\labelenumi{\arabic{enumi}.}
\setcounter{enumi}{15}
\tightlist
\item
  设 E、F 是 K 上的两个向量空间， \(u: E \to F\) 是一个线性映射。若 \(\operatorname{Ker}(u)\) 和 \(\operatorname{Coker}(u)\) 都是有限维的，则称 u 具有有限指标，且数 \(d(u) = \dim(\operatorname{Ker}(u)) - \dim(\operatorname{Coker}(u))\) 称为 u 的指标。
\end{enumerate}

设G是K上的第三个向量空间， \(v: F \rightarrow G\) 一个线性映射。证明：如果三个线性映射 u, v, \(v \circ u\) 中有两个具有有限指标，则第三个也具有有限指标，且 \(d(v \circ u) = d(u) + d(v)\) （利用习题13(a)）。

\begin{enumerate}
\def\labelenumi{\arabic{enumi}.}
\setcounter{enumi}{16}
\tightlist
\item
  设 E, F, G, H 是 K 上的四个向量空间，且 \(u: E \rightarrow F\) , \(v: F \rightarrow G\) , \(w: G \rightarrow H\) 三个线性映射。证明
\end{enumerate}

\[
\operatorname{rg} (v \circ u) + \operatorname{rg} (w \circ v) \leqslant \operatorname{rg} (v) + \operatorname{rg} (w \circ v \circ u).
\]

\begin{enumerate}
\def\labelenumi{\arabic{enumi}.}
\setcounter{enumi}{17}
\item
  设 E、F 是域 K 上的两个向量空间，u、v 是 E 到 F 的两个线性映射。要使存在 E 的自同构 f 和 F 的自同构 g 使得 \(v = g \circ u \circ f\) ，必要且充分的条件是 \(\operatorname{rg}(u) = \operatorname{rg}(v)\) , \(\dim(\operatorname{Ker}(u)) = \dim(\operatorname{Ker}(v))\) 与 \(\dim(\operatorname{Coker}(u)) = \dim(\operatorname{Coker}(v))\) 。
\item
  设 E 是域 K 上的一个向量空间。
\end{enumerate}

\begin{enumerate}
\def\labelenumi{(\alph{enumi})}
\item
  证明每个映射 \(f\) （由 \(\mathbf{E}\) 到 \(\mathbf{E}\) ）只要与 \(\mathbf{E}\) 的所有自同构可交换，就具有形式 \(x \mapsto \alpha x\) ，其中 \(\alpha \in \mathbf{K}\) （首先证明：对所有 \(x \in \mathbf{E}\) 存在 \(\rho(x) \in \mathbf{K}\) 使得 \(f(x) = \rho(x)x\) ，利用 \(f\) 与 \(\mathbf{E}\) 的每个使 \(x\) 不变的自同构可交换这一事实）。
\item
  设 \(g\) 是 \(\mathbf{E} \times \mathbf{E}\) 到 \(\mathbf{E}\) 的一个映射，使得对每个自同构 \(u\) （属于 \(\mathbf{E}\) ）都有 \(g(u(x), u(y)) = u(g(x, y))\) （对所有 \(x, y\) 属于 \(\mathbf{E}\) ）。证明存在两个元素 \(\alpha, \beta\) （属于 \(\mathbf{K}\) ），使得 \(g(x, y) = \alpha x + \beta y\) 在由线性无关元素组成的有序对 \((x, y)\) （这些元素属于 \(\mathbf{E}\) ）所成的集合上成立；此外，存在一个映射 \(\phi\) （由 \(\mathbf{K} \times \mathbf{K}\) 到 \(\mathbf{K}\) ），使得 \(g(\lambda x, \mu x) = \phi(\lambda, \mu)x\) 对所有 \(x \in \mathbf{E}\) 成立（同样的方法）。若进一步 \(g(u(x), u(y)) = u(g(x, y))\) 对每个自同态 \(u\) （属于 \(\mathbf{E}\) ）成立，则 \(g(x, y) = \alpha x + \beta y\) 对所有 \(x, y\) （属于 \(\mathbf{E}\) ）成立。推广到 \(\mathbf{E}^n\) 到 \(\mathbf{E}\) 的映射。
\end{enumerate}

\begin{enumerate}
\def\labelenumi{\arabic{enumi}.}
\setcounter{enumi}{19}
\tightlist
\item
  设 E 是一个向量空间。
\end{enumerate}

\begin{enumerate}
\def\labelenumi{(\alph{enumi})}
\item
  证明：若 M 和 N 是 E 的两个向量子空间，而 M'、N' 分别是 \(\mathbf{M}\) 和 \(\mathbf{N}\) 在 \(\mathbf{E}^*\) 中的正交，则 \(\mathbf{M} \cap \mathbf{N}\) 的正交是 \(\mathbf{M}' + \mathbf{N}'\) （参见 § 2，习题 9 (a)）。
\item
  若 E 是无限维的，证明存在超平面 \(H'\) （属于 \(E^{*}\) ），使得 E 中与 \(H'\) 正交的子空间化为 0（考虑一个包含对应于 E 的一组基的坐标形式的超平面）。
\item
  证明：若 \(\mathbf{E}\) 是无限维的，则存在子空间的一个无限族 \((\mathbf{V}_i)\) （这些子空间属于 \(\mathbf{E}\) ），使得若 \(\mathbf{V}_i'\) 是 \(\mathbf{E}^*\) 中与 \(\mathbf{V}_i\) 正交的子空间，则 \(\mathbf{E}^*\) 中与 \(\bigcap_{i} \mathbf{V}_i\) 正交的子空间不同于 \(\sum_{i} \mathbf{V}_i'\) 。
\item
  由 (b) 推出，若 E 是无限维的，则存在一个作为直和的分解 \(V' \oplus W'\) （这是 \(E^{*}\) 的一个分解），其中 \(W'\) 是有限维的，使得 E 的子空间 V、W 之和 \(V + W\) （其中 V、W 分别与 \(V'\) 和 \(W'\) 正交）不同于 E。
\item
  为使 E 的一个子空间是有限维的，必要且充分条件是它在 \(E^{*}\) 中的正交具有有限余维数。
\end{enumerate}

\begin{enumerate}
\def\labelenumi{\arabic{enumi}.}
\setcounter{enumi}{20}
\item
  用第 6 节定理 8 的记号，设 \(u\) 表示线性映射 \(x \mapsto (\langle x, x_i^* \rangle)\) （由 \(\mathbf{E}\) 到 \(\mathbf{K}_s^{\mathbf{I}}\) ），并设 \(y_0 = (\eta_i)\) 。 \(\mathbf{K}_d^{(\mathbf{I})}\) 与 \(\mathbf{K}_s^{\mathbf{I}}\) 的对偶的一个子空间在 \(\mathbf{K}_d^{(\mathbf{I})}\) 到其双对偶的典范映射下恒等；然后证明：若 \(\mathbf{N}'\) 是 \(^t u\) 的核，则定理 3（第 5 节）的条件表达了下述事实： \(y_0\) 与 \(\mathbf{N}'\) 和 \(\mathbf{K}_d^{(\mathbf{I})}\) 的交正交。当 \(\mathbf{I}\) 无限且方程组 (21)（第 5 节）具有有限秩时，证明这个交不同于 \(\mathbf{N}'\) （注意 \(\mathbf{N}'\) 这时具有有限余维数）。
\item
  设 E 和 F 是域 K 上的两个向量空间，u 是 E 到 F 的一个线性映射。若 V 是 E 的一个向量子空间，而 \(V'\) 是 V 在 \(E^{*}\) 中的正交，证明 \(u(V)\) 的对偶同构于 \({}^{t}u(F^{*})/(V'\cap{}^{t}u(F^{*}))\) 。若 \(W'\) 是 \(F^{*}\) 的一个使 \({}^{t}u(W')\) 为有限维的子空间，而 W 是 \(W'\) 在 F 中的正交，则 \({}^{t}u(W')\) 同构于空间 \(u(E)/(W\cap u(E))\) 的对偶。
\item
  证明：为使线性映射 \(u\) （由向量空间 \(E\) 到向量空间 \(F\) ）满足 \(\mathrm{rg}(t u) = \mathrm{rg}(u)\) ，必要且充分条件是 \(\mathrm{rg}(u)\) 有限（参见习题 3 (d)）。
\item
  设 E 是交换域 K 上有限维数为 n 的向量空间（ \textgreater{} 1）；证明，除非 n = 2 且 K 是两个元素的域，否则不存在仅依赖于 E 的向量空间结构的同构 \(\phi\) （从 E 到 \(E^{*}\) 上）。（注意，若 \(\phi\) 是这样一个同构，则必然对 E 中的 x、y 以及 E 的每个自同构 u 有 \(\langle x, \phi(y) \rangle = \langle u(x), \phi(u(y)) \rangle\) （集合论，IV，§ 1，第 5 节）。）
\item
  \begin{enumerate}
  \def\labelenumii{(\alph{enumii})}
  \tightlist
  \item
    设 K 为一个域，L 与 M 为两个集合；存在典范包含 \((\mathbf{K}_{s}^{\mathbf{L}})^{(\mathbf{M})}\subset(\mathbf{K}_{s}^{(\mathbf{M})})^{\mathbf{L}}\subset\mathbf{K}_{s}^{\mathbf{L}\times\mathbf{M}}\) 。证明：要使这些向量空间中的两个相等，必要且充分条件是集合 L、M 中有一个是有限的。
  \end{enumerate}
\end{enumerate}

\begin{enumerate}
\def\labelenumi{(\alph{enumi})}
\setcounter{enumi}{1}
\tightlist
\item
  设 \((\mathbf{E}_{\lambda})_{\lambda \in L}\) 是 \(K\) 上的一个右向量空间族，而 \((\mathbf{F}_{\mu})_{\mu \in M}\) 是 \(K\) 上的一个左向量空间族。为使典范映射 (26)（第 7 节）是双射，必要且充分条件是向量空间 \(\prod_{\lambda \in L} \mathbf{E}_{\lambda}, \prod_{\mu \in M} F_{\mu}\) 之一是有限维的（利用 (a)）。
\end{enumerate}

\begin{enumerate}
\def\labelenumi{\arabic{enumi}.}
\setcounter{enumi}{25}
\tightlist
\item
  \begin{enumerate}
  \def\labelenumii{(\alph{enumii})}
  \tightlist
  \item
    设 E、F 是域 K 上的两个左向量空间；为使典范同态 \(E^{*} \otimes_{K} F \to \operatorname{Hom}_{K}(E, F)\) 成为双射，必要且充分条件是空间 E、F 之一是有限维的。
  \end{enumerate}
\end{enumerate}

\begin{enumerate}
\def\labelenumi{(\alph{enumi})}
\setcounter{enumi}{1}
\tightlist
\item
  设 E 是一个右向量空间，F 是域 K 上的一个左向量空间。为使典范同态 \(E \otimes_{K} F \to \operatorname{Hom}_{K}(E^{*}, F)\) 成为双射，必要且充分条件是 E 是有限维的（利用习题 3 (d)）。
\end{enumerate}

\begin{enumerate}
\def\labelenumi{\arabic{enumi}.}
\setcounter{enumi}{26}
\tightlist
\item
  \begin{enumerate}
  \def\labelenumii{(\alph{enumii})}
  \tightlist
  \item
    在命题 16 (i)（第 7 节）的假设下，为使典范同态 (27)（第 7 节）成为双射，必要且充分条件是 E 或 F 为有限维（注意 \(\mathrm{Hom}_{\mathbb{K}}(\mathbf{E},\mathbf{G})\otimes_{\mathbb{L}}\mathbf{F}\) 的一个元素在 (27)（第 7 节）下的像把 E 映到 \(\mathbf{G}\otimes_{\mathbb{L}}\mathbf{F}\) 的一个子空间，它包含在形如 \(\mathbf{G}\otimes_{\mathbb{L}}\mathbf{F}'\) 的子空间中，其中 \(\mathbf{F}'\) 是 \(\mathbf{F}\) 的一个有限维子空间）。
  \end{enumerate}
\end{enumerate}

\begin{enumerate}
\def\labelenumi{(\alph{enumi})}
\setcounter{enumi}{1}
\tightlist
\item
  在命题 16 (ii)（第 7 节）的假设下，为使典范同态 (28)（第 7 节）成为双射，必要且充分条件是有序对 \((\mathbf{E}_1, \mathbf{E}_2), (\mathbf{E}_1, \mathbf{F}_1), (\mathbf{E}_2, \mathbf{F}_2)\) 之一由有限维向量空间组成（利用 (a) 及 § 4 的习题 7）。
\end{enumerate}

\begin{enumerate}
\def\labelenumi{\arabic{enumi}.}
\setcounter{enumi}{27}
\item
  设 \(\rho: \mathbf{K} \to \mathbf{A}\) 是域 \(\mathbf{K}\) 到环 \(\mathbf{A}\) 的一个单射同态，并设 \(\mathbf{E}\) 是 \(\mathbf{K}\) 上的一个左向量空间。为使典范映射 \((\mathbf{E}^*)_{(\mathbf{A})} \to (\mathbf{E}_{(\mathbf{A})})^*\) 是双射，必要且充分条件是 \(\mathbf{E}\) 是有限维的，或者 \(\mathbf{A}\) （视为 \(\mathbf{K}\) 上的右向量空间）是有限维的（利用习题 25 ( \(b\) )）。当这两种情形之一成立时，存在典范双射 \((\mathbf{E}^{**})_{(\mathbf{A})} \to (\mathbf{E}_{(\mathbf{A})})^{**}\) 。
\item
  设 A 是一个整环，K 是它的分式域，E 是一个 A-模，而 T(E) 是它的挠模。
\end{enumerate}

\begin{enumerate}
\def\labelenumi{(\alph{enumi})}
\item
  证明 \(\mathbf{E}\) 上的每个线性形式都在 \(\mathrm{T}(\mathrm{E})\) 上为零，从而对偶空间 \(\mathbf{E}^*\) 和 \((\mathrm{E} / \mathrm{T}(\mathrm{E}))^*\) 典范同构。
\item
  证明典范映射 \((\mathbf{E}^{*})_{(\mathbf{K})}\to (\mathbf{E}_{(\mathbf{K})})^{*}\) 是单射，并且当 \(\mathbf{E}\) 有限生成时它是双射。
\item
  给出一个无挠 A-模 E 的例子，使得 \(E^{*} = \{0\}\) 而 \(E_{(K)} \neq \{0\}\) 。
\end{enumerate}

¶ 30. 设 A 是一个整环，E 和 F 是两个 A-模。证明若 x 是 E 的自由元素而 y 是 F 的自由元素，则 \(x \otimes y \neq 0\) （在 \(E \otimes_{A} F\) 中）。（先将其归结为 E 和 F 都无挠的情形，再归结为 E 和 F 都有限生成的情形，并利用习题 29 (b) 证明存在 \(E \times F\) 到 K 的一个 A-双线性映射 f，使得 \(f(x, y) \neq 0\) 。）

*31. 设 \(\mathbf{K}_0\) 为一个交换域，A 为多项式环 \(\mathrm{K}_0[\mathrm{X},\mathrm{Y}]\) （ \(\mathbf{K}_0\) 上两个未定元的多项式环，它是一个整环），而 E 为 A 的理想 \(\mathrm{AX} + \mathrm{AY}\) 。证明在张量积 \(\mathbf{E} \otimes_{\mathbf{A}} \mathbf{E}\) 中，元素 \(\mathbf{X} \otimes \mathbf{Y} - \mathbf{Y} \otimes \mathbf{X}\) 是 \(\neq 0\) 的，并且 \(\mathrm{XY}(\mathbf{X} \otimes \mathbf{Y} - \mathbf{Y} \otimes \mathbf{X}) = 0\) （考虑 \(\mathbf{E} \times \mathbf{E}\) 到商模 \(\mathbf{A}/\mathbf{E}\) 的 A-双线性映射）。*

\begin{enumerate}
\def\labelenumi{\arabic{enumi}.}
\setcounter{enumi}{31}
\item
  将第 10 节的结果推广到 A 为非交换环且具有左分式域 K 的情形（I，§ 9，习题 15）。（注意，若 \(\xi_{i}\) ( \(1 \leqslant i \leqslant n\) ) 是 K 的元素，则存在 A 中的 \(\alpha \neq 0\) ，使得所有 \(\alpha\xi_{i}\) 都属于 A。）
\item
  设 A 是一个整环。A-模 E 称为可除的，如果对所有 \(x \in E\) 以及 A 中所有 \(\alpha \neq 0\) ，存在 \(y \in E\) 使得 \(\alpha y = x\) 。
\end{enumerate}

\begin{enumerate}
\def\labelenumi{(\alph{enumi})}
\item
  设 E、F 是两个 A-模。证明：若 E 是可除的，则 \(E \otimes_{A} F\) 也是可除的。若 E 是可除的且 F 是无挠的，则 \(\operatorname{Hom}_{\mathbf{A}}(\mathbf{E}, \mathbf{F})\) 是无挠的。若 E 是一个挠 A-模而 F 是可除的，则 \(E \otimes_{A} F = \{0\}\) 。若 E 是一个挠 A-模而 F 是无挠的，则 \(\operatorname{Hom}_{\mathbf{A}}(\mathbf{E}, \mathbf{F}) = \{0\}\) 。
\item
  证明每个内射 A-模（§ 2，习题 11）都是可除的；反之，每个可除的无挠 A-模都是内射的（利用 § 2 习题 11 的判据 (δ)）。
\end{enumerate}

\begin{enumerate}
\def\labelenumi{\arabic{enumi}.}
\setcounter{enumi}{33}
\tightlist
\item
  设 A 是一个整环，P 是一个投射 A-模， \(P'\) 是 P 的一个投射子模，而 \(j: P' \to P\) 是典范单射。证明对每个无挠 A-模 E，同态 \(j \otimes 1: P' \otimes_{A} E \to P \otimes_{A} E\) 是单射（归结为 E 有限生成的情形，并将 E 嵌入一个自由 A-模）。
\end{enumerate}

¶ 35. (a) 设 K 为一个域，E 为一个 (K, K)-双模。假设 E 作为 K 上的左向量空间和右向量空间的维数都等于同一个有限数 n。证明存在 E 的 n 个元素的一个族 \((e_{i})\) ，它对于 E 上的两种向量空间结构中的每一种都是基。（注意，若 \((b_{j})_{1\leq j\leq m}\) 是 E 的一个由 m\textless n 个元素组成的族，它对于 E 上的两种向量空间结构中的每一种都是自由的，而 V 是由 \((b_{j})\) 生成的左向量子空间，W 是右向量子空间，则要么 \(V+W\neq E\) ，要么 \(V\cap CW\) 和 \(W\cap CV\) 都非空；在后一情形，若 \(y\in V\cap CW\) , \(z\in W\cap CV\) , 则 \(y+z\) 与 \(b_{j}\) 一起构成 \(m+1\) 个元素的一个族，它对于 E 上的两种向量空间结构中的每一种都是自由的。）

\begin{enumerate}
\def\labelenumi{(\alph{enumi})}
\setcounter{enumi}{1}
\item
  设 F 是 E 的一个子双模，其在 K 上的两种向量空间结构具有相同的维数 p \textless{} n.～证明存在 E 的元素的一个族 \((e_{i})_{1 \leqslant i \leqslant n}\) ，它对于 E 的每一种向量空间结构都是 E 的一组基，并且使得 \((e_{i})_{1 \leqslant i \leqslant p}\) 是 F 在其两种向量空间结构下的基（方法相同）。
\item
  设 \((b_{j})_{1\leqslant j\leqslant n - 1}\) 是 \(n - 1\) 个元素（属于 \(\mathbf{E}\) ）的一个族，它对于 \(\mathbf{E}\) 上的两种向量空间结构中的每一种都是自由的。设 \(\mathbf{V}\) （相应地，W）是 \((b_{j})\) 在 \(\mathbf{E}\) （视为左（相应地，右）向量空间）中生成的超平面。
\end{enumerate}

证明：若 \(\mathbf{V} \subset \mathbf{W}\) （相应地， \(\mathbf{W} \subset \mathbf{V}\) ），则必然 \(\mathbf{V} = \mathbf{W}\) （注意，若 \(a \notin \mathbf{W}\) ，则 \(\lambda \in \mathbf{K}\) 中使 \(\lambda a \in \mathbf{W}\) 的那些 λ 所成的集是一个左理想）。

设 \(\mathbf{K}_0\) 是一个交换域，而 \(\mathbf{K} = \mathbf{K}_0(\mathbf{X})\) 是 \(\mathbf{K}_0\) 上一个未定元的有理函数域（第四章）。在 K 上定义一个 (K, K)-双模结构如下：乘积 \(t.u\) （一个元素 \(u \in \mathbf{K}\) 与一个左算子 \(t \in \mathbf{K}\) 的乘积）是有理函数 \(t(\mathbf{X})u(\mathbf{X})\) ；乘积 \(u.t\) （ \(u\) 与一个右算子 \(t \in \mathbf{K}\) 的乘积）是有理函数 \(u(\mathbf{X})t(\mathbf{X}^2)\) ；这样在 K 上定义的左（相应地，右）向量 K-空间结构是 1 维的（相应地，2 维的）。由此推出 (K, K)-双模的一些例子，使得这种双模的两种向量空间结构的维数是任意整数。

\begin{enumerate}
\def\labelenumi{(\alph{enumi})}
\setcounter{enumi}{1}
\tightlist
\item
  从 (a) 推出一个 (K, K)-双模 E 的例子，其两个向量空间结构具有相同的维数，并且存在 E 的一个子双模 F，它的两个向量空间结构不具有相同的维数。*
\end{enumerate}

设 E 是一个左向量空间（有限维或无限维均可），而 \(A_{i}\) ( \(1 \leqslant i \leqslant n\) ) 是 E 的向量子空间。假设对点的每个序列 \((a_{i})_{1 \leqslant i \leqslant n}\) （这些点属于 E）使得对所有 i 有 \(a_{i} \in A_{i}\) ，由 \(a_{i}\) 生成的 E 的向量子空间都具有维数 \(\leqslant m\) （m 是一个整数 \(\leqslant n\) ）。证明存在 E 的一个维数为 \(h \leqslant m\) 的子空间 W，它包含 \(h + (n - m)\) 个子空间 \(A_{i}\) 。（对 n 作归纳论证。首先证明（对固定的 n）可以只考虑 m \textless{} n 且对所有 i 有 \(\dim(A_{i}) \leqslant m\) 的情形；还可假设对所有 i 有 \(A_{i} \neq \{0\}\) ，且这些 \(A_{i}\) 并不都是

维数为 1 的。然后（对固定的 \(n\) ）对 \(d = \sum_{i=1}^{m} \dim(A_i)\) 作归纳论证。例如，若 \(\dim(A_n) \geqslant 2\) ，考虑一个 1 维子空间 \(B_n\) （包含于 \(A_n\) ），并将归纳假设应用于 \(A_1, \ldots, A_{n-1}, B_n\) ；由此得出结论：存在一个数 \(k\) 使得 \(1 \leqslant k \leqslant m\) ，以及 E 的一个 \(k\) 维子空间 U，它包含 \(k\) 个 \(A_i\) （例如 \(A_1, \ldots, A_k\) ）；如有必要，把 \(k\) 换成一个整数 \(k'\) （满足 \(1 \leqslant k' \leqslant k\) ），并同时证明可以假设存在向量 \(b_i \in A_i\) （对 \(1 \leqslant i \leqslant k\) ）使得 \(b_1, \ldots, b_k\) 构成 U 的一组基；然后投影到 E 中 U 的一个补子空间上，并利用归纳假设。）

\begin{enumerate}
\def\labelenumi{(\alph{enumi})}
\setcounter{enumi}{1}
\tightlist
\item
  设 K 是一个有限域，E 是 K 上的一个 N 维向量空间。设 \(U_{i}\) ( \(1 \leqslant i \leqslant n\) ) 是 E 的向量子空间，使得对 \([1, n]\) 的每个子集 H，诸 \(U_{i}\) （其中 \(i \in H\) ）的交具有维数 \(\leqslant N - \text{Card}(H)\) 。证明这些 \(U_{i}\) 的并不能等于 E（利用对偶性并借助 (a) 论证）。
\end{enumerate}

\begin{enumerate}
\def\labelenumi{\arabic{enumi}.}
\setcounter{enumi}{37}
\item
  设 K 为特征 0 的交换域，E 为 K 上的有限维向量空间。设 \(p_{1}, \ldots, p_{m}\) 为 E 到 E 的向量子空间上的投影算子，使得 \(1_{E} = p_{1} + p_{2} + \cdots + p_{m}\) 。证明 E 是这些 \(p_i(\mathbf{E})\) 的直和，并且这些 \(p_i\) 两两可交换（取两边的迹）。
\item
  设 K 为一个域，L 为 K 的子域，且 \((\mathrm{K}_{\alpha})_{\alpha\in I}\) 为 K 的子域的一个左有向族，使得 \(L=\bigcap_{\alpha\in I}K_{\alpha}\) 。设 E 为 K 上的向量空间，且 \((a_{i})_{1\leqslant i\leqslant m}\) 为 E 的元素的一个有限族；证明若该族 \((a_{i})\) 在 L 上是自由的，则存在一个指标 \(\alpha\) 使得 \((a_{i})\) 在 \(K_{\alpha}\) 上是自由的。（对 m-r 作归纳论证，其中 r 是族 \((a_{i})\) 在 K 上的秩。）
\end{enumerate}

